% !TEX root = ../main.tex
\part{Jogando}

\section{Sequência de Jogo}

O jogo é dividido em turnos, que continuam até que um lado atinja as condições de vitória ou ambos os lados concordem em parar.

\begin{enumerate}[leftmargin=1.2cm]
  \item \textbf{Fase de Sorteio de Dados de Ordem}: certifique-se de que cada unidade não fixada tenha um Dado de Ordem no saco.
  \item \textbf{Fase de Ativação}: os jogadores se revezam, sorteando um Dado de Ordem e ativando a unidade correspondente.
  \item \textbf{Fase de Fim de Turno}: remova todos os Dados de Ordem do campo de batalha (exceto os que estão na posição Abaixar) e prepare-se para o próximo turno.
\end{enumerate}

\section{Cenários e Objetivos}

Os jogos são disputados em cenários que definem a configuração do terreno, a implantação e as condições de vitória.

\subsection{Cenário Introdutório: Batalha de Encontro}

\textbf{Objetivo}: destruir a força inimiga ou causar o máximo de baixas possível.

\textbf{Terreno}: uma área de jogo de 6' x 4' com cerca de 6 a 8 peças de terreno.

\textbf{Implantação}: as forças se posicionam nas bordas opostas da mesa (60cm de distância).

\subsection{Ideia de Cenário: O Caminho Seguro}

Um cenário que exercita Engenheiros e cooperação de fogo:

\textbf{Mapa}: 3 linhas de arame farpado e 2 campos minados no meio da mesa.

\textbf{Objetivo}: a Squad Atacante precisa atravessar o mapa e colocar uma Carga de Demolição no Bunker Defensor no final da mesa.

\textbf{Desafio}: o Defensor tem uma LMG no Bunker; o Atacante precisa usar Engenheiros para cortar o arame enquanto o resto do esquadrão dá cobertura contra a LMG.

\begin{nota}
Cenários adicionais e tabelas de missões detalhadas podem ser expandidos em um apêndice futuro.
\end{nota}
